% Chapter 26: Sunflowers, Thresholds, and the Erdős-Faber-Lovász Conjecture
\let\LocalMacro\relax

\chapter{Sunflowers, Thresholds, and the Erd\H{o}s-Faber-Lov\'{a}sz Conjecture}

\section{The Problem}

\subsection{Informal}

\textbf{Background.} A sunflower is a collection of groups of objects that all overlap in exactly the same common core, with each group contributing unique ``petals'' outside the core. For example, three groups $\{1,2,3\}$, $\{1,2,4\}$, and $\{1,2,5\}$ form a sunflower with core $\{1,2\}$ and petals $\{3\}$, $\{4\}$, $\{5\}$.

\textbf{Given.} You are given a large collection of groups, each containing exactly $w$ objects, drawn from a shared pool of objects.

\textbf{Goal.} Determine how many groups you need before a sunflower with a fixed number of petals is guaranteed to appear, regardless of how the groups are chosen.

\textbf{Constraints.} Each group has exactly $w$ elements. The sunflower must have a common core (possibly empty) shared by all petals. The bound must hold for all possible collections, with no exceptions. The question is about the minimum number of groups needed to force a sunflower.

\textbf{Example.} If each group has 3 elements and you want a sunflower with 3 petals, how many groups do you need? With 4 groups, you might avoid sunflowers. The sunflower conjecture predicts that for each fixed number of petals $k$, a family of more than $c_k^w$ distinct $w$-element sets must contain a sunflower, where $c_k$ depends only on $k$.

\subsection{Formal}
\begin{conjecture}[Erd\H{o}s-Rado, 1960]
For each fixed $k$, there is a constant $c_k$ such that every family of more than $c_k^w$ distinct sets of size $w$ contains a sunflower with $k$ petals.
\end{conjecture}

Erd\H{o}s offered \$1000 for a proof. The best bound before 2022 was $(w!)^{k-1}$.

\subsection{Historical Context}
Three major conjectures in this area were resolved or substantially advanced around 2021--2022, but by different methods. Alweiss, Lovett, Wu, and Zhang improved the sunflower bound using spread approximations; Park and Pham proved the Kahn--Kalai expectation-threshold conjecture; and Kang, Kelly, K\"uhn, Methuku, and Osthus proved the Erd\H{o}s--Faber--Lov\'{a}sz conjecture using methods including absorption.

\subsection{What Was Known}
The sunflower conjecture was open since 1960, with the best bound growing factorially. The Kahn--Kalai conjecture had been open since 2006, though Kalai had earlier suggested the threshold should be close to the expectation. The Erd\H{o}s--Faber--Lov\'{a}sz conjecture had been open since 1972, with only partial results for small values of $n$.

\subsection{Why It Matters}
These three problems are fundamental questions about how structured patterns emerge from large collections of objects. Sunflowers appear in data analysis and computational biology. Threshold phenomena govern how networks become connected. Graph coloring is central to scheduling, register allocation in compilers, and frequency assignment in wireless networks.

\subsection{Simplified Version}
For the sunflower problem, when the number of petals is fixed at $k = 3$, the Erd\H{o}s--Rado bound is $(w!)^2$. The conjecture asks for a bound of the form $c_k^w$. Alweiss--Lovett--Wu--Zhang improved the known bound to roughly $(C_k\log w)^w$ for fixed $k$.

\subsection{The Sunflower Conjecture}
A \emph{sunflower} with $k$ petals is a family of $k$ sets whose pairwise intersections are all equal to a common core $C$.

\subsection{The Kahn-Kalai Conjecture (Expectation-Threshold)}
For an increasing property $\mathcal{P}$ on a finite ground set $X$, its threshold $p_c(\mathcal{P})$ is the value of $p$ at which a random subset belongs to $\mathcal{P}$ with probability $1/2$. The expectation threshold $q(\mathcal{P})$ is defined using weighted covers of the property, as described in the Mathematical Theory section; it is not simply the first-moment threshold for counting all minimal witnesses.

\begin{conjecture}[Kahn-Kalai, 2006]
$p_c(\mathcal{P}) \le C \log \ell \cdot q(\mathcal{P})$, where $\ell$ is the maximum size of a minimal witness (at least 2).
\end{conjecture}

\section{Mathematical Theory}

\subsection{Prerequisites}
This chapter assumes familiarity with:
\begin{itemize}
\item Combinatorics (set families, extremal combinatorics)
\item Basic probability (expectations, concentration inequalities)
\item Graph theory (chromatic number, cliques)
\end{itemize}

\subsection{Sunflowers and the Erd\H{o}s-Rado Bound}

A \emph{sunflower} with $k$ petals and core $C$ is a family $\mathcal{F} = \{A_1, \ldots, A_k\}$ of sets such that $A_i \cap A_j = C$ for all $i \neq j$. The sets $A_i \setminus C$ are pairwise disjoint (the ``petals'').

\begin{definition}[Sunflower]
A family $\{A_1, \ldots, A_k\}$ of subsets of $[w]$ is a \emph{sunflower} if there exists a set $C \subset [w]$ (the \emph{core}) such that $A_i \cap A_j = C$ for all $i \neq j$.
\end{definition}

The Erd\H{o}s--Rado bound states that any family of $w$-element sets with more than $(w!)^{k-1}$ members contains a $k$-petal sunflower. For fixed $k$, the conjecture asks for a bound of the form $c_k^w$ instead.

\subsection{Spread Families}

The key innovation of Alweiss--Lovett--Wu--Zhang is the notion of a \emph{spread} family:

\begin{definition}[Spread Family]
A family $\mathcal{F}$ of subsets of $[w]$ is \emph{$s$-spread} if every element $i \in [w]$ appears in at most $|\mathcal{F}| / w^s$ sets of $\mathcal{F}$.
\end{definition}

\textbf{Concrete example.} Suppose $w = 100$ and $|\mathcal{F}| = 10{,}000$. If $\mathcal{F}$ is 1-spread, each element $i \in [100]$ appears in at most $10{,}000 / 100^1 = 100$ sets. If $\mathcal{F}$ is 2-spread, each element appears in at most $10{,}000 / 100^2 = 1$ set---so no element appears in more than one set at all. The parameter $s$ controls how ``uniform'' the distribution is: larger $s$ means more uniform. The power $w^s$ in the denominator means that in a spread family, no single element can appear in too many sets relative to the total---the ``frequency'' of each element is at most $|\mathcal{F}|/w^s$, which is small when $s$ is large. The idea is that in a spread family, no element ``dominates''---every element appears with roughly equal frequency. Any family can be ``spread-approximated'' by partitioning it into subfamilies and discarding sets containing over-represented elements.

\subsection{Thresholds and the Kahn-Kalai Framework}

For an increasing property $\mathcal{P}$ on the Boolean lattice $2^{[n]}$, the \emph{threshold} is:
\[
p_c(\mathcal{P}) = \inf\{p : \Pr_{G \sim G(n,p)}[G \in \mathcal{P}] \ge 1/2\}.
\]

The expectation threshold is defined using weighted covers, rather than just counting all minimal witnesses. A family $\mathcal{G}$ of subsets of the ground set is a $p$-small cover if every set in $\mathcal{P}$ contains some $S\in\mathcal{G}$ and $\sum_{S\in\mathcal{G}}p^{|S|}\leq 1/2$. The expectation threshold $q(\mathcal{P})$ is the largest $p$ for which such a cover exists. This definition accounts for smaller subsets that already force the property.

The Kahn--Kalai theorem bounds the threshold in terms of this expectation threshold:

\begin{fact}
$q(\mathcal{P}) \le p_c(\mathcal{P}) \le C \log \ell \cdot q(\mathcal{P})$, where $\ell$ is the maximum size of a minimal witness (at least 2) and $C$ is an absolute constant.
\end{fact}

\subsection{The Absorption Method}

The absorption method is a technique for completing partial constructions. The idea is:

\begin{enumerate}
\item Reserve a small ``absorbing set'' $A$ at the start.
\item Build the desired structure on the remaining vertices, possibly leaving a small remainder $R$.
\item Use the absorbing set to ``absorb'' $R$: the set $A$ is designed so that it can accommodate any small remainder.
\end{enumerate}

The power of absorption is that it reduces a global problem to a local one: instead of building the entire structure at once, you build most of it and then clean up.

\section{The Solution}

\subsection{The Big Idea}

For the sunflower problem, Alweiss, Lovett, Wu, and Zhang used spread approximations: they reduce the search to families in which no element is too heavily represented, then use probabilistic and entropy arguments to find a sunflower. The Kahn--Kalai and Erd\H{o}s--Faber--Lov\'{a}sz conjectures were resolved by separate arguments and should not be presented as consequences of this technique.

\subsection{What Was Stopping Progress}

The Erd\H{o}s--Rado sunflower conjecture (1960) asks how many $w$-element sets guarantee a sunflower with $k$ petals. The classical bound is roughly $(w!)^{k-1}$, while the conjecture predicts $c_k^w$ for fixed $k$. The challenge is to remove the factorial dependence on set size without assuming special structure in the family.

\subsection{The Breakthrough}

Alweiss, Lovett, Wu, and Zhang (published in 2021) introduced \textbf{spread approximations}: a technique for reducing arbitrary set families to ``spread'' families where no element is too heavily represented. Their argument combines this reduction with probabilistic estimates to improve the classical sunflower bound substantially.

\subsection{How It Works}

The breakthrough of Alweiss, Lovett, Wu, and Zhang (2022) \cite{alweiss2022} was the introduction of \textbf{spread approximations}---a technique for reducing arbitrary set families to ``spread'' families where every element appears roughly uniformly.

\begin{definition}[Spread Family]
A family $\mathcal{F}$ of subsets of $[w]$ is \emph{$s$-spread} if every element $i \in [w]$ appears in at most $|\mathcal{F}| / w^s$ sets of $\mathcal{F}$.
\end{definition}

The key insight is twofold:

\begin{enumerate}
\item \textbf{Spread approximation.} The reduction extracts structured subfamilies with controlled element frequencies. The quantitative loss depends on the parameters; it is not a single logarithmic factor independent of the sunflower bound.
\item \textbf{Spread families are easy.} In a spread family, no element dominates. A random subset of a spread family of the right size contains a sunflower with high probability, because the uniformity prevents any single element from ``blocking'' all potential sunflowers. The proof uses entropy arguments and the Lopsided Lov\'{a}sz Local Lemma.
\end{enumerate}

For each fixed number of petals $k$, their result gives a bound of order $(C_k\log w)^w$ on the number of $w$-element sets needed to force a sunflower. This improves the classical factorial-scale bound, though it does not prove the conjectured $c_k^w$ bound.

\begin{theorem}[ALWZ \cite{alweiss2022}]
For each fixed $k$, any sufficiently large family of $w$-element sets with more than $(C_k\log w)^w$ members contains a sunflower with $k$ petals, for a constant $C_k$ depending only on $k$.
\end{theorem}

Their proof uses:
\begin{enumerate}
\item \textbf{Spread approximations}: A set family is ``spread'' if no element appears too often. Any family can be approximated by a spread family.
\item \textbf{Random sampling}: A random subset of a spread family contains a sunflower with high probability.
\item \textbf{Entropy arguments}: Use the Lopsided Lov\'{a}sz Local Lemma to find the sunflower.
\end{enumerate}

\begin{highlightbox}
Spread approximation is a useful tool in extremal combinatorics, but the Kahn--Kalai and Erd\H{o}s--Faber--Lov\'{a}sz proofs use different frameworks.
\end{highlightbox}

\subsection{Sunflowers: Alweiss, Lovett, Wu, Zhang (2022)}

\begin{theorem}[ALWZ \cite{alweiss2022}]
Any family of sets of size $w$ with more than $(C \log w)^{k-1} w!$ sets contains a sunflower with $k$ petals.
\end{theorem}

Their proof uses:
\begin{enumerate}
\item \textbf{Spread approximations}: A set family is ``spread'' if no element appears too often. Any family can be approximated by a spread family.
\item \textbf{Random sampling}: A random subset of a spread family contains a sunflower with high probability.
\item \textbf{Entropy arguments}: Use the Lopsided Lov\'{a}sz Local Lemma to find the sunflower.
\end{enumerate}

\subsection{Kahn-Kalai: Park and Pham (2022)}

\begin{theorem}[Park-Pham \cite{parkpham2022}]
The Kahn-Kalai conjecture is true: $p_c(\mathcal{P}) \le C \log \ell \cdot q(\mathcal{P})$.
\end{theorem}

Their proof uses spread approximations combined with weighted random sampling and entropy compression.

\subsection{Erd\H{o}s-Faber-Lov\'{a}sz: Kang, Kelly, K\"{u}hn, Methuku, Osthus (2022)}

\begin{theorem}[KKKMO \cite{kang2022}]
The EFL conjecture is true for sufficiently large $n$.
\end{theorem}

Their proof uses the \emph{absorption method}: build the coloring iteratively, absorbing leftover vertices at the end, using expansion properties of the hypergraph to guarantee absorption works.

\begin{highlightbox}
Spread approximation is a useful tool in extremal combinatorics, but it was not the common method behind all three results.
\end{highlightbox}

\subsection{After the Proof}

The Kahn--Kalai expectation-threshold conjecture and the Erd\H{o}s--Faber--Lov\'{a}sz conjecture were proved in 2022. The sunflower conjecture remains open: the Alweiss--Lovett--Wu--Zhang result substantially improved the known bound but did not reach the conjectured $c_k^w$ form.

The spread approximation method has become a new standard technique alongside the polynomial method and the probabilistic method. It has been applied to other problems in combinatorics where uniform distributions play a role, and its connections to entropy methods and the Lov\'{a}sz Local Lemma have been further explored. The resolution of the Kahn--Kalai conjecture, in particular, unified two branches of probabilistic combinatorics---random graph thresholds and expected thresholds---and the proof techniques continue to influence the field.

For sunflowers, the known bound of roughly $(C_k\log w)^w$ remains larger than the conjectured $c_k^w$. The gap is an open problem. The Kahn--Kalai theorem is proved, although sharper constants and related threshold questions continue to be studied.

\section{Impact}
\begin{itemize}
\item \textbf{Extremal combinatorics}: Spread approximations form a new toolkit alongside the polynomial method and the probabilistic method.
\item \textbf{Threshold phenomena}: The Kahn-Kalai result connects random graph thresholds to expected thresholds, unifying two branches of probabilistic combinatorics.
\item \textbf{Graph coloring}: The EFL result settles a 49-year-old conjecture about the chromatic number of structured hypergraph unions.
\end{itemize}

\section{Computational Verification}

The sunflower problem and related threshold phenomena can be explored computationally by searching for sunflower-free families and verifying the spread approximation technique.

\subsection{Finding Sunflower-Free Families}

\begin{verbatim}
from itertools import combinations

def is_sunflower(sets, k=3):
    """Check if any k sets form a sunflower."""
    n = len(sets)
    for combo in combinations(range(n), k):
        # Check if sets[i] intersect in the same core
        intersection = sets[combo[0]]
        for s in combo[1:]:
            intersection = intersection & s
        # Check all pairwise intersections equal the core
        is_sf = True
        for i in range(k):
            for j in range(i+1, k):
                if sets[combo[i]] & sets[combo[j]] != intersection:
                    is_sf = False
                    break
            if not is_sf:
                break
        if is_sf:
            return True, combo, intersection
    return False, None, None

def find_sunflower_free(w, max_sets=200):
    """Greedy search for sunflower-free family of w-element sets."""
    from itertools import combinations as combs
    universe = set(range(w))
    family = []
    for r in range(w, w + max_sets):
        for s in combs(universe, r):
            candidate = family + [set(s)]
            sf, _, _ = is_sunflower(candidate)
            if not sf:
                family.append(set(s))
                if len(family) >= max_sets:
                    break
    return family

w = 5
family = find_sunflower_free(w, 50)
sf, combo, core = is_sunflower(family)
print(f"Sunflower-free family of {w}-element sets: {len(family)} sets")
print(f"Contains sunflower: {sf}")
\end{verbatim}

\subsection{Spread Approximation Demo}

\begin{verbatim}
def compute_element_frequencies(family, w):
    """Count how often each element appears in the family."""
    freq = [0] * w
    for s in family:
        for elem in s:
            freq[elem] += 1
    return freq

def make_spread(family, w, s):
    """Remove sets containing over-represented elements."""
    n = len(family)
    threshold = n / (w ** s)
    spread_family = []
    for s_set in family:
        if all(compute_element_frequencies([s_set], w)[e] <= threshold
               for e in s_set):
            spread_family.append(s_set)
    return spread_family

w = 10
family = find_sunflower_free(w, 30)
freq = compute_element_frequencies(family, w)
print(f"Element frequencies: {freq}")
spread = make_spread(family, w, 0.5)
print(f"After spread approximation: {len(spread)} sets remain")
\end{verbatim}

\subsection{Threshold Phenomena in Random Graphs}

\begin{verbatim}
import random

def has_triangle(adj, n):
    """Check if graph has a triangle."""
    for i in range(n):
        for j in range(i+1, n):
            if adj[i][j]:
                for k in range(j+1, n):
                    if adj[i][k] and adj[j][k]:
                        return True
    return False

def random_graph_threshold(n, trials=1000):
    """Estimate the triangle-free threshold for G(n, p)."""
    results = {}
    for p_int in range(5, 95, 5):
        p = p_int / 100
        count = 0
        for _ in range(trials):
            adj = [[False]*n for _ in range(n)]
            for i in range(n):
                for j in range(i+1, n):
                    if random.random() < p:
                        adj[i][j] = adj[j][i] = True
            if has_triangle(adj, n):
                count += 1
        results[p] = count / trials
    return results

results = random_graph_threshold(20, 500)
for p, frac in sorted(results.items()):
    if 0.2 < frac < 0.8:
        print(f"p={p:.2f}: P(has triangle) = {frac:.2f}")
\end{verbatim}

\begin{highlightbox}
Students should run these scripts to observe: (1)~sunflower-free families can be quite large, (2)~the spread approximation removes over-represented elements, and (3)~random graph thresholds exhibit sharp phase transitions.
\end{highlightbox}

\section{Exercises}

\begin{exercise}[Basic]
Show that a sunflower with $k$ petals and core $C$ has the property that any two petals intersect exactly in $C$. Find a sunflower with 3 petals in $\{1,2,3\}, \{1,4,5\}, \{1,6,7\}$.
\end{exercise}

\begin{exercise}[Intermediate]
Explain the difference between the threshold $p_c(\mathcal{P})$ and the expectation threshold $q(\mathcal{P})$. Why is $q(\mathcal{P}) \le p_c(\mathcal{P})$ always?
\end{exercise}

\begin{exercise}[Challenge]
Describe the absorption method intuitively: why does ``building in spare capacity'' allow one to complete a coloring at the end?
\end{exercise}

\section{Timeline}

\begin{tabularx}{\linewidth}{lX}
\toprule
\textbf{Year} & \textbf{Event} \\ \midrule
1960 & Erd\H{o}s-Rado sunflower conjecture \\
1972 & Erd\H{o}s-Faber-Lov\'{a}sz conjecture \\
2006 & Kahn-Kalai conjecture on expectation-threshold \\
2022 & ALWZ prove sunflower conjecture (up to log factor) \\
2022 & Park-Pham prove Kahn-Kalai conjecture \\
2022 & KKKMO prove EFL conjecture \\
2022 & Huh wins Fields Medal (log-concavity in combinatorics) \\
\bottomrule
\end{tabularx}

\section{Citations}

\begin{enumerate}
\item Alweiss, R., Lovett, S., Wu, K., \& Zhang, J. (2022). ``Improved bounds for the sunflower lemma.'' Annals of Mathematics, 196(3), 977--1021.
\item Park, J., \& Pham, H. T. (2022). ``A proof of the Kahn-Kalai conjecture.'' arXiv:2203.17207.
\item Kang, D., Kelly, T., K\"{u}hn, D., Methuku, A., \& Osthus, D. (2022). ``A proof of the Erd\H{o}s-Faber-Lov\'{a}sz conjecture.'' arXiv:2101.04686.
\item Erd\H{o}s, P., \& Rado, R. (1960). ``Combinatorial theorem on classifications of sets of elements.'' Acta Mathematica Hungarica, 11, 117--142.
\item Kahn, J., \& Kalai, G. (2006). ``A counterexample to Borsuk's conjecture.'' Bulletin of the AMS, 29, 60--68.
\end{enumerate}
